\scanpage{166}
% archive foliation: 165
\deleted{$\unsure{\Pi.\ S}$} $P_i$ se déduit de $P_i'$ par une \uline{section} de $P_i'/\pi$ passant par $\unsure{\bar\xi'|\bar\xi_i'}$. $S$ étant connexe, il ne peut y en avoir qu'une seule telle section.

c) Si $\pi\mathbin{-}\pi'\rightrightarrows\pi''$ est exacte, alors \deleted{$\unsure{\Phi}$} $\Phi(\pi)\mathbin{-}\Phi(\pi')\rightrightarrows\Phi(\pi'')$ aussi.

Il \deleted{\unsure{reste}} à prouver que si $P'$ de groupe $\pi'$, pt marqué $\xi'$, est tel que \deleted{\unsure{on}} l'on ait un isom
\[
P'\times_{\pi'}(\pi'',\varphi_1)\xrightarrow[\sim]{u}
P'\times_{\pi'}(\pi'',\varphi_2),
\]
alors on peut restreindre le groupe structural.

Or soit $Q$ le sous-schéma de $P'$ des coïncidences de $u i_1$ et $i_2$. Il est ouvert et fermé, contient $\xi'$, est stable par $\pi$, donc correspond à un sous-schéma de $P'/\pi$ à la fois ouvert et fermé $S'$ [donc \inserted{absolument} fidèlement plat et quasi-compact].

\margin{\diagram{\begin{tikzcd}
\pi \arrow[r,dash] & \pi' \arrow[r,shift left=.5ex] \arrow[r,shift right=.5ex] & \pi''
\end{tikzcd}}
\diagram{\begin{tikzcd}
P' \arrow[r,"i_1"] \arrow[dr,"i_2"'] & P_1'' \arrow[d,"u"]\\
& P_2''
\end{tikzcd}}}

De plus, on voit que les \deleted{\unsure{deux}} fibres géom de $S'\mathbin{-}S$ ont \unsure{partout} un pt; donc c'est un isom $S'\to S$ \unsure{alors}, \unsure{ce qui prouve}\ldots
